\section{轨道控制与 Monte Carlo 仿真}
\label{sec:control}

不稳定轨道（Halo/NRHO/DPO）必须配合轨道保持才能真正可用。e2m2e 把
控制仿真作为设计链的最后一环：设计阶段产出的 NominalOrbit（等间隔状态
表 + Floquet 基 + 投影因子）直接被控制算法消费，Floquet 分析结果无需
重算。

\subsection{三种控制律}

\begin{itemize}
  \item \textbf{特征点控制}：只在轨道的特征位置（如近月点、平面穿越点）
        施加修正，修正次数少、燃耗低，适合长期保持；
  \item \textbf{严格目标点控制}：把状态精确打回标称轨道上的目标点，
        跟踪精度最高；
  \item \textbf{松弛目标点控制}：目标点容许误差带，在跟踪精度与推进
        消耗之间折中。
\end{itemize}

此外提供角动量管理功能：姿态与推力器联合控制，把动量卸载纳入保持策略。

\subsection{Monte Carlo 误差仿真}

轨道保持的燃耗统计必须在误差模型下评估。e2m2e 支持导航误差与推力误差
联合注入的 Monte Carlo 仿真：对同一标称轨道批量运行保持仿真，统计
$\Delta V$ 消耗分布与轨道维持情况，为任务论证给出置信区间而非单点
估值。仿真产物自动归档入轨道库并支持标签管理——在 MCP 场景下，
“对这条 NRHO 跑 100 次 Monte Carlo 保持仿真并打上
\texttt{candidate} 标签”是一条自然语言指令即可完成的工作流。
